\uuid{WI0V}
\titre{Application linéaire/Nombres complexes}
\auteur{hanani ; mzari}
\organisation{exo7}
\licence{CC BY-NC-SA 4.0}
\datecreate{2020-01-20}
\niveau{L1}
\module{Algèbre}
\chapitre{Application linéaire}
\sousChapitre{Définition}
\difficulte{2}
\theme{Applications linéaires, Application linéaire, Nombres complexes}
\exo7qcmid{lille-2.applications-lineaires.005}

\contenu{
\begin{question}
\qclassification{107.01}
\qtags{motcle=Application linéaire/Nombres complexes}
On considère les applications suivantes :
$$\begin{array}{rccc}f:&\Cc&\to&\Cc\\
& z&\to& \Re (z)\end{array}   \quad \mbox{et} \quad  \begin{array}{rccc}g:&\Cc&\to&\Cc\\
& z&\to& \Im (z), \end{array}$$
où $\Re (z)$ (resp. $\Im (z)$) est la partie réelle (resp. imaginaire) de $z$. Quelles sont les assertions vraies ?
\begin{answers}
\bad{$f$ est $\Cc$-linéaire}.
\good{$f$ est $\Rr$-linéaire}.
\good{$g$ est  $\Rr$-linéaire}.
\bad{$g$ est $\Cc$-linéaire}.
\end{answers}
\begin{explanations} On vérifie que $f$ et $g$ sont $\Rr$-linéaires. Par contre, elles ne sont pas $\Cc$-linéaires.
\end{explanations}
\end{question}
}
